\documentclass[11pt]{article}
\usepackage[margin=1in]{geometry}
\usepackage{amsmath,amssymb,amsthm}
\usepackage[hidelinks]{hyperref}
\newtheorem{theorem}{Theorem}
\newtheorem{lemma}{Lemma}
\newcommand{\CT}{\operatorname{CT}}
\newcommand{\bn}[2]{\binom{#1}{#2}}
\title{The constant-term identity in OEIS A220119}
\author{Shreyas Samir Kolte}
\date{November 23, 2026}
\begin{document}
\maketitle
\begin{abstract}
We give a finite coefficient-extraction proof of the Laurent-polynomial
identity proposed by F.~Chapoton in OEIS A220119 on May 25, 2026.
The argument uses the binomial theorem, Vandermonde's identity, and a
factorial identity. It holds for every nonnegative integer $n$.
The separate divisibility conjecture in that entry is not addressed.
\end{abstract}

Write $\bn uv=0$ when $v<0$ or $v>u$, for $u\geq0$.
All variables below are commuting indeterminates, and all coefficient
extractions and sums are finite.
Put
\[
 F_n(k)=\bn nk^2\bn{n+k}n,\qquad
 A_n=\sum_{j,k=0}^{n}F_n(j)F_n(k)\bn{j+k}n.
\]
The double sum is the definition used in OEIS A220119; it also appears
in Almkvist and Zudilin~\cite[Section 4]{AZ}.

\begin{theorem}\label{thm:main}
For every integer $n\geq0$,
\[
 A_n=\CT_{a,b,c,d,e}
 \left(
 \frac{(1+c)(1+c+d)(1+a)(1+a+c+d+e)
 (1+a+b)(1+a+b+c+d+e)}
 {abcde}
 \right)^n.
\]
\end{theorem}

\section*{An auxiliary coefficient identity}
For $0\leq\ell\leq n$, define
\[
 H_{n,\ell}=[c^nd^n](1+c)^n(1+c+d)^n(c+d)^\ell.
\]

\begin{lemma}\label{lem:aux}
For $0\leq\ell\leq n$,
\begin{equation}\label{eq:aux}
 \bn{2n-\ell}n H_{n,\ell}
   =\sum_{k=0}^{n}F_n(k)\bn{k}{n-\ell}.
\end{equation}
\end{lemma}
\begin{proof}
Expanding in the total degree in $c,d$ gives
\begin{equation}\label{eq:Hfirst}
 H_{n,\ell}=\sum_{r=0}^{\ell}
 \bn nr\bn n{\ell-r}\bn{n+r}n
 =[x^\ell](1+x)^n L_n(x),
\end{equation}
where
\[
 L_n(x)=\sum_{r=0}^{n}\bn nr\bn{n+r}n x^r.
\]
Indeed, in the expansion of $(1+c)^n$, a term $c^{n-r}$
must be paired with $(c+d)^{n+r}$ to contribute to $c^nd^n$.

The following polynomial identity is a direct consequence of Vandermonde:
\begin{equation}\label{eq:L}
 L_n(x)=\sum_{r=0}^{n}\bn nr^2 x^r(1+x)^{n-r}.
\end{equation}
To verify it, the coefficient of $x^t$ on the right is
\[
 \sum_{r=0}^{t}\bn nr^2\bn{n-r}{t-r}
 =\bn nt\sum_{r=0}^{t}\bn nr\bn tr
 =\bn nt\bn{n+t}n.
\]
Combining \eqref{eq:Hfirst} and \eqref{eq:L} yields
\begin{equation}\label{eq:Hsecond}
 H_{n,\ell}=\sum_{r=0}^{\ell}
 \bn nr^2\bn{2n-r}{\ell-r}.
\end{equation}
For $0\leq r\leq\ell\leq n$, cancellation of factorials gives
\begin{equation}\label{eq:factorial}
 \bn{2n-\ell}n\bn nr^2\bn{2n-r}{\ell-r}
 =\bn n\ell\bn nr\bn\ell r\bn{2n-r}n.
\end{equation}
Both sides of \eqref{eq:factorial} equal
\[
 \frac{n!(2n-r)!}
 {(r!)^2((n-r)!)^2(n-\ell)!(\ell-r)!}.
\]
Consequently, the left side of \eqref{eq:aux} is
\[
 \bn n\ell\sum_{r=0}^{\ell}\bn nr\bn\ell r\bn{2n-r}n.
\]
For its right side, use
$\bn nk\bn{k}{n-\ell}
 =\bn n\ell\bn\ell{n-k}$ and then set $r=n-k$:
\[
 \sum_{k=0}^{n}\bn nk^2\bn{n+k}n\bn{k}{n-\ell}
 =\bn n\ell\sum_{r=0}^{\ell}
   \bn nr\bn\ell r\bn{2n-r}n.
\]
These expressions agree.
\end{proof}

\section*{Proof of the theorem}
\begin{proof}[Proof of Theorem \ref{thm:main}]
Let $B_n$ denote the asserted constant term. Set
$X=1+a$ and $T=1+a+c+d=X+c+d$.
Equivalently, $B_n$ is the coefficient of $a^nb^nc^nd^ne^n$ in
\[
 (1+c)^n(1+c+d)^n X^n(X+b)^n(T+e)^n(T+b+e)^n.
\]
First,
\[
 [e^n](T+e)^n(T+b+e)^n
 =\sum_{i=0}^{n}\bn ni^2 T^{n-i}(T+b)^i.
\]
Next,
\[
 [b^n](X+b)^n(T+b)^i
 =\sum_{\substack{0\leq j\leq n\\n-j\leq i}}
     \bn nj\bn i{n-j}X^{n-j}T^{i-n+j}.
\]
After multiplication by $X^nT^{n-i}$, the monomial in $X,T$
is $X^{2n-j}T^j$. Its weight, summed over $i$, is
\begin{equation}\label{eq:weight}
 \bn nj\sum_{i=0}^{n}\bn ni^2\bn i{n-j}
 =\bn nj^2\bn{n+j}n=F_n(j).
\end{equation}
For completeness, the first equality follows from
$\bn ni\bn i{n-j}=\bn nj\bn j{n-i}$ and
$\sum_i\bn ni\bn j{n-i}=\bn{n+j}n$.
Thus
\[
 B_n=\sum_{j=0}^{n}F_n(j)R_{n,j},
\]
where
\[
 R_{n,j}=[a^nc^nd^n](1+c)^n(1+c+d)^n
                    X^{2n-j}(X+c+d)^j.
\]
Expand the last factor in powers of $c+d$ and extract $a^n$:
\[
 R_{n,j}=\sum_{\ell=0}^{j}
 \bn j\ell\bn{2n-\ell}n H_{n,\ell}.
\]
Lemma \ref{lem:aux}, followed by Vandermonde's identity, gives
\begin{align*}
 R_{n,j}
 &=\sum_{k=0}^{n}F_n(k)
       \sum_{\ell=0}^{j}\bn j\ell\bn k{n-\ell}\\
 &=\sum_{k=0}^{n}F_n(k)\bn{j+k}n.
\end{align*}
Substitution proves $B_n=A_n$. Every displayed identity also holds
when $n=0$, so the proof includes that case.
\end{proof}

\section*{Verification and status}
The proof above is an all-$n$ algebraic argument, not an inference from
initial values. The accompanying standard-library Python script independently
expands the five-variable numerator for $0\leq n\leq10$ and checks the
auxiliary identities for $0\leq n\leq45$, using exact integers.
The accompanying formalization passed the recorded Lean 4.19.0 build and audit. Independent human peer review is not claimed.

The accessible OEIS entry retrieved on November 23, 2026 still labels this
specific constant-term identity as a conjecture and attributes it to
F.~Chapoton, May 25, 2026. Targeted searches did not locate an earlier proof;
this is not a guarantee of priority. The classical identities used in the
argument are not claimed to be new. The separate claim
$(n+1)(n+2)\mid A_n$ for $n>0$ is not proved here.

\begin{thebibliography}{9}
\bibitem{OEIS}
OEIS Foundation Inc.,
\emph{A220119}, especially the comment by F.~Chapoton dated May 25, 2026,
\url{https://oeis.org/A220119}.

\bibitem{AZ}
G.~Almkvist and W.~Zudilin,
\emph{Differential equations, mirror maps and zeta values},
arXiv:math/0402386, version 2 (2004), Section 4.
\url{https://arxiv.org/abs/math/0402386}.
\end{thebibliography}
\paragraph{Assistance.} Substantial OpenAI ChatGPT/Codex assistance was used in the proof, writing, formalization, and debugging.
\end{document}
